\section{Introduction}
\label{sec:introduction}

This paper is a certificate paper for a fixed geometric object, not a general
hinged-dissection theorem.  The object is the median-plane S4 dissection of a
regular tetrahedron.  The four pieces are tetrahedra; at the closed endpoint
they tile the original tetrahedron, and along two audited one-parameter signed
rays they are moved by rooted rigid hinge-tree kinematics.

The certificate is useful precisely because the tempting stronger statements
are false or unproved.  The package does not prove that a positive-thickness
printable mechanism exists.  It does not prove that every possible S4 hinge
choice works.  It does not prove dynamic connectedness between the two audited
representative signed rays.  It proves a scoped statement for two representative
trees and a one-dimensional angular domain, with every pair obligation routed
to a recorded certificate layer.

\paragraph{Contribution.}
The package contributes four pieces of evidence.
\begin{enumerate}[leftmargin=*]
  \item An exact coordinate model for the S4 median-plane dissection and its
  closed-contact endpoint.
  \item A rooted rigid-hinge kinematic model for \treeA{} and \treeB{} with
  explicit hinge axes and signed-ray conventions.
  \item A route-wrapper theorem on the domain
  \(\theta=0\) plus \(0<\theta\le 120^\circ\), supported by the A6--A7d
  certified layers.
  \item A local B04/A7c bridge that upgrades six selected-hinge signed rows
  from mere orientation signs to contact-side/no-strict-local-overlap records.
\end{enumerate}

Replay-backed claims are labelled as \emph{Certified Propositions}.  These are
finite or symbolic certificate statements tied to files in the repository, not
claims of a new general theorem in the literature.

\paragraph{Public boundary.}
The claim boundary is part of the result.  The package allows the zero-thickness
one-parameter route-wrapper claim and the six-row selected-hinge bridge claim.
It disallows physical validation, positive-thickness hinge validity, global
collision-free motion, and three-parameter bounded-cell closure.  These
boundaries are mirrored in \path{docs/PUBLIC_CLAIM_BOUNDARY.md},
\path{docs/CLAIM_LEDGER.md}, and \path{docs/PROOF_OBLIGATIONS.md}.

\paragraph{Organization.}
Section~\ref{sec:model} gives the exact tetrahedral model.  Section~\ref{sec:endpoint}
records the closed endpoint and kinematic conventions.  Section~\ref{sec:wrapper}
states the scoped one-parameter theorem wrapper.  Section~\ref{sec:layers}
explains the A6--A7d certificate layers.  Section~\ref{sec:b04} proves the
selected-hinge wedge bridge used by the B04 rows.  Sections~\ref{sec:package}
and~\ref{sec:rw12} describe reproducibility and the RW12 digital-fabrication
boundary.  Section~\ref{sec:related} places the result next to related hinged
mechanism and fabrication work.
